% ============================================================================
%  98-thuat-ngu.tex — bảng thuật ngữ Việt–Anh cho toàn cây
%  Ngày tạo: 2026-09-07 | Sửa gần nhất: 2026-09-07
% ============================================================================
\documentclass[english.tex]{subfiles}
\begin{document}
\chapter*{Bảng thuật ngữ}
\ptrtarget{Z/98}
\addcontentsline{toc}{chapter}{Bảng thuật ngữ}
\markboth{Bảng thuật ngữ}{}

\begin{tomtat}
Bảng này gom các thuật ngữ dùng xuyên suốt quyển sách, kèm dạng tiếng Anh và một câu giải thích. Cột cuối trỏ tới chương nói kỹ nhất về khái niệm đó. Dạng tiếng Anh ở đây quan trọng hơn bình thường: đó là dạng bạn sẽ gặp khi tra cứu và khi đọc tài liệu về ngôn ngữ, và nhiều thuật ngữ không có bản dịch tiếng Việt đã ổn định --- đúng hiện tượng mà \ptr{B/10} mô tả.
\end{tomtat}

\section*{A. Nghĩa và từ vựng}
\begin{center}\footnotesize
\rowcolors{2}{white}{tblalt}
\begin{longtable}{@{}L{3.4cm}L{3.8cm}L{6.6cm}L{1.5cm}@{}}
\toprule
\textbf{Tiếng Việt} & \textbf{English} & \textbf{Nghĩa ngắn} & \textbf{Chương}\\
\midrule
\endfirsthead
\toprule
\textbf{Tiếng Việt} & \textbf{English} & \textbf{Nghĩa ngắn} & \textbf{Chương}\\
\midrule
\endhead
nghĩa nằm ở cách dùng & meaning is use & Nghĩa của một từ được xác định bởi các ngữ cảnh nó xuất hiện & \ptr{A/01}\\
độ phủ từ vựng & lexical coverage & Tỉ lệ từ trong văn bản mà bạn đã biết & \ptr{A/02}\\
vốn từ nhận ra & receptive vocabulary & Các từ bạn hiểu khi gặp & \ptr{A/02}\\
vốn từ chủ động & productive vocabulary & Các từ bạn dùng được khi viết và nói & \ptr{E/19}\\
từ vựng học thuật chung & academic vocabulary & Từ dùng nhiều trong mọi ngành: \emph{derive}, \emph{constitute} & \ptr{A/02}\\
từ kỹ thuật & technical word & Từ chỉ tồn tại hoặc chỉ mang nghĩa đó trong một ngành & \ptr{B/10}\\
từ bán kỹ thuật & semi-technical word & Từ thông thường mang nghĩa hẹp trong ngành --- loại nguy hiểm nhất & \ptr{B/10}\\
họ từ & word family & Một từ gốc cùng mọi dạng biến đổi và phái sinh của nó & \ptr{A/04}\\
gốc, tiền tố, hậu tố & root, prefix, suffix & Ba thành phần cấu tạo từ & \ptr{A/04}\\
văn vực & register & Mức trang trọng và bối cảnh dùng của một từ & \ptr{D/16}\\
màu sắc đánh giá & semantic prosody & Thái độ khen hay chê gắn liền với từ, dù nghĩa mô tả trung tính & \ptr{D/16}\\
gần nghĩa & near-synonym & Hai từ nghĩa cốt lõi gần nhau nhưng khác văn vực hoặc phạm vi & \ptr{D/16}\\
kết hợp từ & collocation & Các từ quen đi cùng nhau hơn mức ngẫu nhiên & \ptr{D/17}\\
cụm đóng gói & formulaic sequence & Chuỗi từ được lưu và dùng như một đơn vị & \ptr{D/17}\\
lỗi chuyển di & transfer error & Lỗi do dịch thẳng tổ hợp từ tiếng mẹ đẻ & \ptr{D/17}\\
ngữ liệu & corpus & Kho văn bản thật dùng để tra cách dùng thực tế & \ptr{D/16}\\
\bottomrule
\end{longtable}
\end{center}

\section*{B. Đọc và hiểu ý}
\begin{center}\footnotesize
\rowcolors{2}{white}{tblalt}
\begin{longtable}{@{}L{3.4cm}L{3.8cm}L{6.6cm}L{1.5cm}@{}}
\toprule
\textbf{Tiếng Việt} & \textbf{English} & \textbf{Nghĩa ngắn} & \textbf{Chương}\\
\midrule
\endfirsthead
\toprule
\textbf{Tiếng Việt} & \textbf{English} & \textbf{Nghĩa ngắn} & \textbf{Chương}\\
\midrule
\endhead
ba lượt đọc & three-pass reading & Đọc một bài báo theo ba lượt với ba mục đích khác nhau & \ptr{B/07}\\
đọc rộng & extensive reading & Đọc nhiều tài liệu dễ để củng cố; nhánh 1 & \ptr{E/19}\\
khẳng định và bằng chứng & claim and evidence & Hai thành phần cơ bản của một lập luận & \ptr{B/06}\\
cơ sở suy luận & warrant & Giả định nối bằng chứng với khẳng định & \ptr{B/06}\\
siêu diễn ngôn & metadiscourse & Ngôn ngữ nói về chính văn bản và về thái độ tác giả & \ptr{B/08}\\
rào đón & hedging & Ngôn ngữ hạ mức cam kết: \emph{may}, \emph{suggest} & \ptr{C/14}\\
tăng cường & booster & Ngôn ngữ nâng mức cam kết: \emph{clearly}, \emph{demonstrate} & \ptr{C/14}\\
mức cam kết & epistemic stance & Độ mạnh mà một câu khẳng định điều gì đó là đúng & \ptr{C/14}\\
hàm ý & implicature & Điều người nói ngụ ý mà không nói ra & \ptr{B/08}\\
định nghĩa thao tác & operational definition & Định nghĩa bằng cách nói đo thế nào & \ptr{B/10}\\
định nghĩa theo loại và khác biệt & genus and differentia & Nêu loại lớn hơn rồi nêu đặc điểm phân biệt & \ptr{B/10}\\
tín hiệu cấu trúc & discourse signal & Cụm báo hiệu chuyển ý trong bài giảng hoặc bài viết & \ptr{B/09}\\
\bottomrule
\end{longtable}
\end{center}

\section*{C. Viết}
\begin{center}\footnotesize
\rowcolors{2}{white}{tblalt}
\begin{longtable}{@{}L{3.4cm}L{3.8cm}L{6.6cm}L{1.5cm}@{}}
\toprule
\textbf{Tiếng Việt} & \textbf{English} & \textbf{Nghĩa ngắn} & \textbf{Chương}\\
\midrule
\endfirsthead
\toprule
\textbf{Tiếng Việt} & \textbf{English} & \textbf{Nghĩa ngắn} & \textbf{Chương}\\
\midrule
\endhead
danh từ hoá & nominalization & Biến động từ hoặc tính từ thành danh từ & \ptr{C/11}\\
động từ rỗng & light verb & Động từ không mang hành động thật: \emph{make}, \emph{perform} & \ptr{C/11}\\
chủ thể và hành động & character and action & Nguyên tắc: chủ thể làm chủ ngữ, hành động làm động từ & \ptr{C/11}\\
vị trí chủ đề & topic position & Đầu câu --- nơi người đọc tìm mối nối & \ptr{C/12}\\
vị trí nhấn & stress position & Cuối câu --- nơi trọng âm rơi và trí nhớ giữ lại & \ptr{C/12}\\
cũ trước mới sau & given--new & Bắt đầu câu bằng thông tin đã biết, kết thúc bằng thông tin mới & \ptr{C/12}\\
móc xích & chaining & Phần mới cuối câu này thành phần cũ đầu câu sau & \ptr{C/12}\\
câu chủ đề & topic sentence & Câu mở đoạn nêu ý chính của đoạn & \ptr{C/13}\\
IMRAD & IMRAD & Khuôn bốn phần của bài báo thực nghiệm & \ptr{C/13}\\
ba nước đi mở bài & CARS moves & Dựng bối cảnh, chỉ khoảng trống, chiếm chỗ trống & \ptr{C/13}\\
khoảng trống & research gap & Chỗ mà kiến thức hiện có chưa xử lý & \ptr{C/13}\\
phạm vi & scope & Giới hạn điều kiện mà một khẳng định đúng & \ptr{C/14}\\
quy trách & attribution & Nói rõ một khẳng định là của ai & \ptr{C/14}\\
lời nguyền tri thức & curse of knowledge & Không thấy được chỗ khó hiểu vì bạn đã biết ý mình & \ptr{C/15}\\
sửa theo lượt & pass-based revision & Mỗi lượt chỉ tìm một loại vấn đề & \ptr{C/15}\\
\bottomrule
\end{longtable}
\end{center}

\section*{D. Ngôn từ và luyện tập}
\begin{center}\footnotesize
\rowcolors{2}{white}{tblalt}
\begin{longtable}{@{}L{3.4cm}L{3.8cm}L{6.6cm}L{1.5cm}@{}}
\toprule
\textbf{Tiếng Việt} & \textbf{English} & \textbf{Nghĩa ngắn} & \textbf{Chương}\\
\midrule
\endfirsthead
\toprule
\textbf{Tiếng Việt} & \textbf{English} & \textbf{Nghĩa ngắn} & \textbf{Chương}\\
\midrule
\endhead
ẩn dụ ý niệm & conceptual metaphor & Hiểu một miền trừu tượng qua một miền cụ thể & \ptr{D/18}\\
miền nguồn, miền đích & source, target domain & Miền được mượn và miền cần hiểu & \ptr{D/18}\\
ẩn dụ chết & dead metaphor & Ẩn dụ dùng nhiều tới mức không ai còn thấy là ẩn dụ & \ptr{D/18}\\
định khung & framing & Chọn cách mô tả khiến người nghe nghĩ theo một hướng & \ptr{D/18}\\
bốn nhánh & four strands & Tiếp nhận dễ, học chủ đích, sản xuất, luyện lưu loát & \ptr{E/19}\\
luyện lưu loát & fluency development & Làm việc dễ hơn khả năng, có áp lực tốc độ & \ptr{E/19}\\
luyện có chủ đích & deliberate practice & Nhắm điểm yếu cụ thể, khó vừa quá sức, có phản hồi & \ptr{E/19}\\
ôn tập giãn cách & spaced repetition & Giãn dần khoảng cách giữa các lần ôn & \ptr{E/20}\\
hiệu ứng kiểm tra & testing effect & Tự nhớ lại củng cố trí nhớ mạnh hơn đọc lại & \ptr{E/19}\\
mức tối thiểu & minimum dose & Lượng nhỏ nhất bạn làm kể cả ngày bận nhất & \ptr{E/21}\\
\bottomrule
\end{longtable}
\end{center}
\end{document}
